\chapter{Introducción}\label{cha:introduccion}

\section{Presentación del Documento}

El presente informe describe el proyecto de desarrollo Gestión de repositorios semánticos compatible con el estándar \acrshort{oaipmh}, un aplicativo que pretende extender a \acrshort{labman}, el \acrlong{dms} de los grupos de Internet y Telecomunicaciones de DeustoTech, detallando tanto los objetivos que se pretenden alcanzar con el proyecto, como las fases, actividades y recursos necesarios para llevarlo a cabo.

El contenido de este documento se estructura en torno a los siguientes productos:

\begin{itemize}
	\item \textbf{Definición de proyecto:}
		
	Establecimiento del objetivo fundamental del proyecto, especificando cuáles son los aspectos funcionales que lo comprenden y cuáles son los que quedan excluidos.
	
	\item \textbf{Producto final:}
		
	Especificación de la solución elegida que va a construir el proyecto en cuestión.
	
	\item \textbf{Descripción de la realización:}

	Realización y definición de las diferentes actividades cuyo desarrollo va a permitir la realización y consecución del objetivo del proyecto.

	\item \textbf{Organización:}
	
	Definición del equipo de trabajo que desarrollará el proyecto, así como su estructura organizativa, sistema de gestión y seguimiento del trabajo.

	\item \textbf{Condiciones de ejecución:}

	Definición del entorno de trabajo, de los criterios sobre los que se van a realizar las sucesivas recepciones, así como el tratamiento que se va a establecer para aquellos casos que puedan ser considerados como modificaciones o mejoras en el planteamiento inicial del proyecto.

	\item \textbf{Planificación:}
	Estimación de cargas y duración de las diferentes actividades del proyecto, así como su asignación a los diferentes miembros del equipo y su planificación en el tiempo.

	\item \textbf{Valoración económica:}
	Determinación del valor correspondiente a este proyecto, de los hitos de facturación y de la forma de pago.
\end{itemize}

\section{Introducción a LabMan}

Es este apartado se realiza una breve introducción a \acrfull{labman}, el \acrlong{dms} de MoreLab, el grupo de investigación formado por los equipos de Internet y Telecomunicaciones de DeustoTech y en el que se basa este proyecto.

\begin{figure}[!htp]
	\centering
	\includegraphics[scale=0.15]{fig/morelab-logo}
	\caption{Logo de MORELab}
\end{figure}

Gestionar la información no siempre es una tarea trivial y más aún cuando hay que tratar con los datos que componen varias entidades como pueden ser los proyectos, investigadores, publicaciones, eventos, etc. dentro del ámbito de la investigación. La mayoría de los grupos de investigación utilizan sistemas de gestión de contenido tales como Joomla!\cite{joomla}, WordPress\cite{wordpress} o Drupal\cite{drupal} para exponer sus datos. Sin embargo, para extraer la información de estos \acrshortpl{cms} se requieren herramientas externas para llevar a cabo técnicas de análisis de datos. 
Para hacer uso de estas herramientas normalmente hace falta generar documentos adicionales, tales como \acrshort{csv}, ficheros de texto, entre otros que provocan redundancia de la información, que provoca dificultades a la hora de actualizar los datos y la calidad de los mismos. Esta situación empeora cuando además se disponen de distintas fuentes para la obtención de información, como pueden ser las paginas web personales de los investigadores en los que se muestran sus logros a lo largo de su carrera, la información financiera gestionada por su propio departamento, etc.

Del esfuerzo para gestionar la información grupo de investigación de MoreLab nace {labman}, con el objetivo de gestionar todo este tipo información diferenciadose de otros \acrshort{cms} por apostar por la exposición de los datos como \acrlong{lod}\cite{linkeddata}.

Los datos enlazados\cite{pena_visual_2014}
\section{Introducción a SQL}

Como su propio nombre indica, \acrfull{sql} es un lenguaje de programación, estandarizado por \acrshort{iso}\cite{ISO} y \acrshort{ansi}\cite{ANSI}, diseñado para gestionar la información almacenada en los sistemas de gestión de bases de datos relacionales, como por ejemplo PostgreSQL\cite{PostgreSQL}, MySQL\cite{MySQL}, SQLite\cite{SQLite}, por medio de consultas estructuradas en inglés.

Las consultas representan las operaciones más comunes y esenciales del lenguaje \acrshort{sql} para la recopilación de información dentro de una \acrshort{bd}. Estas consultas se realizan por medio de la sentencia ``SELECT'', que pueden ser complementadas por medio de clausulas para realizar búsquedas específicas. De todas las clausulas disponibles hay que destacar la siguientes:

\begin{itemize}
	\item \textbf{FROM:} Indica de que tabla o tablas de debe extraerse la información.
	\item \textbf{WHERE:} Sirve para delimitar las tuplas o filas de la tablas en las que se realiza la búsqueda. Si las filas no cumplen las condiciones expuestas en esta clausula, serán excluidas del resultado.
	\item \textbf{ORDER BY:} Esta es la única forma para definir el criterio de ordenación los resultados en \acrshort{sql}, sin esta clausula el orden sería aleatorio.
\end{itemize}

\lstinputlisting[language=SQL, frame=single, label={lst:selectsql}, caption=Ejemplo de sentencia SELECT en \acrshort{sql}]{content/code/sql/select-example.sql}

En esta consulta \acrshort{sql} (ver algoritmo \ref{lst:selectsql}) se solicitan todos los nombres y las ciudades de los clientes que residan en Suecia.
\section{Introducción a SPARQL}

\acrshort{sparql}\cite{SPARQL_language} es el acrónimo recursivo del inglés \acrlong{sparql} que hace referencia tanto a el lenguaje estandarizado por la \acrshort{rdf} \acrshort{dawg} de la \acrshort{w3c}\cite{W3C} para consultas a grafos \acrshort{rdf} como para el protocolo de invocación de consultas \acrshort{sparql} remotas.

El lenguaje de consultas \acrshort{sparql} (actualmente en la versión 1.1), permite tanto buscar como manipular grafos \acrshort{rdf} disponibles en la web o bases de datos semánticas almacenados como tripletas, en otras palabras, define un lenguaje equivalente a \acrshort{sql} a excepción de que este se utiliza exclusivamente para bases de datos semánticas.

Al igual que en \acrshort{sql} las consultas en \acrshort{sparql} constituyen las operaciones más comunes y esenciales del lenguaje. Su estructura es es muy similar a la ya vista en el algoritmo \ref{lst:selectsql}, haciendo uso de la misma sentencia ``SELECT''.

\lstinputlisting[language=SQL, otherkeywords={PREFIX}, frame=single, label={lst:selectsparql}, caption=Ejemplo de sentencia SELECT en \acrshort{sparql}]{content/code/sparql/select-example.sparql}

Esta consulta \acrshort{sparql} (ver algoritmo \ref{lst:selectsparql}) añade un nuevo elemento ``PREFIX'' a lo anteriormente visto en \acrshort{sql}, que tiene como función almacenar \acrshortpl{uri} para reducir la longitud de de las mismas a la hora de acceder a sus atributos. La consulta en sí busca en los sujetos ?planttype y los objetos ?name que estén relacionado con el predicativo plant:planttype y devuelve solo sus nombres.